\begin{frame}{Estrategia de Búsqueda y Optimización}

    \begin{columns}[T]
        %========================================
        % COLUMNA IZQUIERDA: Algoritmos
        %========================================
        \column{0.45\textwidth}
        
        \begin{block}{1. Nelder-Mead (Principal)}
            \scriptsize
            \begin{itemize}
                \item Algoritmo numérico iterativo.
                \item Busca el mínimo de la función de error sin usar derivadas.
            \end{itemize}
        \end{block}

        \vspace{0.3cm}

        \begin{block}{2. Búsqueda por Grilla (Soporte)}
            \scriptsize
            \begin{itemize}
                \item Exploración espacial estructurada.
                \item Útil para visualización conceptual, no como optimizador final.
            \end{itemize}
        \end{block}

        %========================================
        % COLUMNA DERECHA: Rangos
        %========================================
        \column{0.50\textwidth}
        
        \begin{block}{Acotación del Dominio ($x, y, z$)}
            \centering
            \scriptsize
            \vspace{0.1cm}
            \begin{tabular}{@{}lcc@{}}
                \toprule
                \textbf{Eje} & \textbf{Rango} & \textbf{Justificación Física} \\ \midrule
                $x$ & $[-5, 5]$ & Cobertura horizontal de la red. \\ \vspace{0.1cm}
                $y$ & $[-5, 5]$ & Cobertura horizontal de la red. \\ \vspace{0.1cm}
                $z$ & $[-3, 1]$ & Profundidad del hipocentro. \\ \bottomrule
            \end{tabular}
        \end{block}

        \vspace{0.4cm}
        \begin{center}
            \footnotesize
            \textbf{Limitar el dominio asegura la convergencia numérica.}
        \end{center}
    \end{columns}

\end{frame}